\section{Related Work}

% \jiabin{It's great that you have all the relevant papes listed. Can you provide a STORYLINE in each topic as bullet-points?
% Within each topic, it will be good to have two paragraphs: 1) What people have done? and 2) How are those papers relate to our work?
% }

% \devi{Might be too risky; but I wonder if we can say that we already discussed related work in intro, and for a more detailed discussion of existing relevant work, see Appendix? We could say this somewhere in the intro?}

In this section, we discuss the related work in video generation using different models. 
% including Generative Adversarial Networks (GAN), Autoregressive models (AR), and implicit neural representations (INR). 
We focus on the different strategies adopted by these methods to handle temporal dynamics and their potential for and challenges associated with long video generation. Also see supp. mat. D for other relevant models and tasks.

\topic{GAN-based video generator.} Adapting GANs~\cite{karras2019style,karras2020analyzing,goodfellow2014generative} to video synthesis requires modeling the temporal dimension. Both 3D deconvolutionals~\cite{vondrick2016generating} and additional RNN or LSTM~\cite{saito2017temporal,Tulyakov_2018_CVPR,Tulyakov_2018_CVPR,clark2019adversarial,saito2020train,munoz2021temporal,tian2021a} have been used. By unrolling the steps taken by the RNNs, videos of longer duration can be generated. However, as shown in our experiments, the quality of these videos degrades quickly. 
% The length of videos generated by these models is limited by the GPU memory~\cite{castrejon2021hierarchical}
In addition, without further modifications, the length of videos generated by these models is limited by the GPU memory (e.g., at most 140 frames can be generated on 32GB V100 GPU by HVG~\cite{castrejon2021hierarchical}).

\topic{AR-based video generator.} Autoregressive models have become a ubiquitous generative model for video synthesis~\cite{ranzato2014video,srivastava2015unsupervised,kalchbrenner2017video,Weissenborn2020Scaling}. 
% \cite{ranzato2014video,srivastava2015unsupervised} are the first to utilize RNN and LSTM for video prediction. Video Pixel Network~\cite{kalchbrenner2017video} generalizes  PixelCNN~\cite{oord2016conditional} for video generation. More recently, Video Transformer~\cite{Weissenborn2020Scaling} builds its architecture on the Transformer~\cite{vaswani2017attention}. 
A common challenge faced by AR models is their slow inference speed. This issue is mitigated by training on the compressed tokens with VQVAEs~\cite{rakhimov2020latent,wu2021n,le2021ccvs,yan2021videogpt}. 
% Furthermore, NÜWA~\cite{wu2021n} unifies the training of generative models for images and videos using a frame-based VQGAN and adopts 3D nearby self-attention for better efficiency. 
% A frame-based compression often leads to temporal artifacts and does not fully utilize temporal redundancy. To this end, CCVS~\cite{le2021ccvs} utilizes a flow module to improve temporal consistency. VideoGPT~\cite{yan2021videogpt} further leverages 3D convolution and axial self-attention. 
Our model falls into this line of video generators. We show that such a VQVAE-based AR model is promising to generate long videos with long-range dependence.

\topic{INR-based video generator.} Implicit Neural Representations~\cite{sitzmann2020implicit,tancik2020fourier} represent continuous signals such as videos by mapping the coordinate space to RGB space~\cite{yu2021generating,skorokhodov2021stylegan}. 
% Some recent works apply it to image synthesis~\cite{skorokhodov2021adversarial,anokhin2021image} by training generators to generate the weights of the neural networks. 
% DIGAN~\cite{yu2021generating} first generalized this idea to video generation by decomposing the network weights for the spatial and temporal coordinates. StyleGAN-v~\cite{skorokhodov2021stylegan} maps a continuous temporal positional embedding to the input feature map of the StyleGAN. 
The advantage of these models is their ability to generate arbitrarily long videos non-autoregressively. However, their generated videos still suffer from the quality degradation or contain periodic artifacts due to the positional encoding and struggle at synthesizing new content. 
% \devi{Do we show this in our results? If so, we should say ``as we empirically demonstrate''}.

\topic{Concurrent works on long video generation.} \cameraready{In parallel to our work, \cite{ho2022video} proposes a gradient conditioning method for sampling longer videos with a diffusion model, \cite{nash2022transframer} and  \cite{hong2022cogvideo} explore a hierarchical model with frame-level VQVAE embeddings or RGB frames. \cite{brooks2022generating} uses a low-resolution long video generator and short-video super-resolution network to generate videos of dynamic scenes.} 
% \devi{I made minor edits; please check for accuracy. What are ``scene videos''?}

% \devi{Xi's literature review was on text to video generation. Do we want to include a paragraph on that? Several of the papers Xi found can generate arbitrary length. Some of them are on just moving digits and not realistic video.}

% \xiyin{I shared the survey video/slides to Songwei. If we have text-conditioned video generation experiments (I assumed we have), it might be worth adding a section about conditional video generation. One of the paper did claim they can generate arbitrary length, that is the advantage of the early LSTM based approaches. and as Devi mentioned, most of them are evaluated on very constrained datasets like moving digits, human walking, and simple UCF actions.}
% \jiabin{I think it will be good to have a separate topic on PADDING. }